\chapter{多项式逼近}
\label{chap:polynomial-approximation}

Taylor 多项式来自函数在一点的导数，只在余项可控时才逼近原函数。连续函数
可能处处不可导，却仍能在整个闭区间上被多项式一致逼近。这里的多项式不再
由某一点的局部数据决定，而是由函数在许多采样点的值共同构造。

本章先给出 Bernstein 的显式构造，再证明更一般的 Stone--Weierstrass
定理。前者回答“怎样逼近”，后者说明“哪些函数代数足以逼近”。

\section{稠密性与一致逼近}

设 $A$ 是赋范空间 $X$ 的子集。若对每个 $x\in X$ 和每个
$\varepsilon>0$，都存在 $a\in A$ 使 $\norm{x-a}<\varepsilon$，则称 $A$
在 $X$ 中稠密，等价地 $\overline A=X$。

对 $C([a,b])$ 的多项式子空间，稠密性意味着：任给连续函数 $f$ 与误差
$\varepsilon>0$，可以找到一个多项式 $p$，使
\[
 |f(x)-p(x)|<\varepsilon\qquad(a\le x\le b)
\]
同时成立。这里的多项式次数可以依赖于 $f$ 和 $\varepsilon$，但不能依赖于
$x$。逐点为每个 $x$ 另选一个多项式并不能得到一致逼近。

稠密子空间不必等于整个空间。多项式集合只是可数维线性空间，而
$C([a,b])$ 是无限维的；极限过程填补了二者之间的差距。若限制多项式次数
不超过固定 $N$，所得有限维子空间是闭的，不可能稠密于 $C([a,b])$。

\section{Bernstein 多项式}

先在 $[0,1]$ 上工作。对 $f\in C([0,1])$，定义第 $n$ 个 Bernstein 多项式
\[
 B_nf(x)=\sum_{k=0}^n f\left(\frac kn\right)
 \binom nk x^k(1-x)^{n-k}.
\]
基函数非负，且由二项式定理
\[
 \sum_{k=0}^n\binom nkx^k(1-x)^{n-k}=1.
\]
因此 $B_nf(x)$ 是采样值 $f(k/n)$ 的凸组合，特别有
$\norm{B_nf}_\infty\le\norm f_\infty$，并保持常数和函数值的上下界。
映射 $B_n:C([0,1])\to C([0,1])$ 是正线性算子：$f\ge0$ 推出
$B_nf\ge0$，并且 $B_n1=1$。这些性质使误差估计可以直接从函数值之差
出发，而无须控制多项式系数。

三个矩恒等式决定了逼近误差：
\begin{align}
 B_n1(x)&=1,\\
 B_n(t)(x)&=x,\\
 B_n((t-x)^2)(x)&=\frac{x(1-x)}n.
 \label{eq:bernstein-moments}
\end{align}
第二式可对 $(u+v)^n$ 求导后令 $u=x,v=1-x$ 得到；再求一次导数可算出
$k(k-1)$ 的和，从而得到第三式。

定义连续模
\[
 \omega_f(\delta)=\sup\{|f(s)-f(t)|:s,t\in[0,1],\ |s-t|\le\delta\}.
\]
一致连续性说明 $\omega_f(\delta)\to0$。若 $|s-t|>\delta$，把区间分成
不超过 $1+|s-t|/\delta$ 段，便有
\begin{equation}
 |f(s)-f(t)|\le
 \left(1+\frac{|s-t|}{\delta}\right)\omega_f(\delta).
 \label{eq:modulus-telescoping}
\end{equation}

\begin{theorem}[Bernstein 逼近]
\label{thm:bernstein-uniform-approximation}
对每个 $f\in C([0,1])$，$B_nf\to f$ 一致成立。更具体地，
\[
 \norm{B_nf-f}_\infty
 \le\omega_f(\delta)
 \left(1+\frac{1}{2\delta\sqrt n}\right)
 \qquad(\delta>0).
\]
\end{theorem}

\begin{proof}
记 $p_{n,k}(x)=\binom nkx^k(1-x)^{n-k}$。由
\eqref{eq:modulus-telescoping}，
\[
 |B_nf(x)-f(x)|
 \le\omega_f(\delta)
 \left(1+\frac1\delta\sum_{k=0}^n
 \left|\frac kn-x\right|p_{n,k}(x)\right).
\]
对有限和使用 Cauchy--Schwarz，并利用 \eqref{eq:bernstein-moments}，
\[
 \sum_k\left|\frac kn-x\right|p_{n,k}(x)
 \le\left(\sum_k\left(\frac kn-x\right)^2p_{n,k}(x)\right)^{1/2}
 \le\frac1{2\sqrt n}.
\]
估计与 $x$ 无关。取 $\delta=n^{-1/4}$，右端趋于零。
\end{proof}

这个证明没有假设 $f$ 可导。逼近来自两点：权重集中在 $k/n\approx x$ 的
位置，连续函数在相近点的值相近。若 $f$ 是 Lipschitz 函数，直接取
$|f(k/n)-f(x)|\le L|k/n-x|$ 可得
$\norm{B_nf-f}_\infty\le L/(2\sqrt n)$。

\section{Weierstrass 逼近定理}

\begin{theorem}[Weierstrass]
\label{thm:weierstrass-polynomial-approximation}
实多项式在 $C([a,b];\R)$ 的上确界范数下稠密。
\end{theorem}

\begin{proof}
仿射变换 $x=a+(b-a)t$ 把 $[0,1]$ 映到 $[a,b]$。令
$g(t)=f(a+(b-a)t)$。定理 \ref{thm:bernstein-uniform-approximation} 给出
多项式 $B_ng(t)$ 一致逼近 $g$。把 $t=(x-a)/(b-a)$ 代回，仍得到 $x$ 的
多项式，并一致逼近 $f$。若 $a=b$，常多项式已经足够。
\end{proof}

逼近多项式一般不是 Taylor 多项式。连续函数没有导数时，Taylor 展开甚至
无从定义；即使函数无穷次可导，Taylor 级数也可能不等于原函数。Weierstrass
定理只要求每个误差水平存在某个多项式，并不要求这些多项式由同一组局部
导数生成。

Bernstein 算子还保持单调性和凸性。记
$\Delta f_k=f((k+1)/n)-f(k/n)$，直接求导得
\[
 (B_nf)'(x)
 =n\sum_{k=0}^{n-1}\Delta f_k
 \binom{n-1}{k}x^k(1-x)^{n-1-k}.
\]
若 $f$ 单调递增，则每个 $\Delta f_k\ge0$，所以 $B_nf$ 也单调递增。
再求一次导数，
\[
 (B_nf)''(x)
 =n(n-1)\sum_{k=0}^{n-2}
 \left[f\!\left(\frac{k+2}{n}\right)
 -2f\!\left(\frac{k+1}{n}\right)
 +f\!\left(\frac{k}{n}\right)\right]
 \binom{n-2}{k}x^k(1-x)^{n-2-k}.
\]
凸函数在等距节点上的二阶差分非负，故 $B_nf$ 保持凸性。这些保形性质是
显式构造相对于纯存在性证明的优点；Weierstrass 定理本身只给存在性，不会
自动给出这样的形状控制。

\section{分离点的函数代数}

设 $K$ 为紧致 Hausdorff 空间，$A\subset C(K;\R)$ 是子代数。称 $A$
含常数，如果所有常函数属于 $A$；称 $A$ 分离 $K$ 中的点，如果任意
$x\ne y$，都存在 $f\in A$ 使 $f(x)\ne f(y)$。

分离点是必要条件：若 $f(x)=f(y)$ 对每个 $f\in A$ 都成立，那么该等式在
$\overline A$ 中仍成立，$A$ 不可能逼近一个在 $x,y$ 取不同值的连续函数。
含常数则保证能同时调节函数的平移和尺度。

证明一般定理还需要一个格闭包事实。

\begin{lemma}[代数闭包的格性质]
\label{lem:stone-weierstrass-lattice-closure}
若 $A\subset C(K;\R)$ 是含常数的子代数，则其一致闭包
$\overline A$ 对绝对值、有限最大值和有限最小值封闭。
\end{lemma}

\begin{proof}
给定 $f\in\overline A$，设 $M=\norm f_\infty$。由 Weierstrass 定理，在
$[-M,M]$ 上存在实多项式 $p_n$ 一致逼近 $|t|$；可把常数项稍作调整使
$p_n(0)=0$。由于 $\overline A$ 是闭子代数且含常数，$p_n\circ f\in
\overline A$，并一致趋于 $|f|$，故 $|f|\in\overline A$。再利用
\[
 f\vee g=(f+g+|f-g|)/2,
 \qquad f\wedge g=(f+g-|f-g|)/2
\]
得到有限最大、最小值的封闭性。
\end{proof}

\section{Stone--Weierstrass 定理}

\begin{theorem}[实 Stone--Weierstrass 定理]
\label{thm:real-stone-weierstrass}
设 $K$ 为紧致 Hausdorff 空间。若子代数 $A\subset C(K;\R)$ 含常数并
分离点，则 $A$ 在 $C(K;\R)$ 中一致稠密。
\end{theorem}

\begin{proof}
固定 $f\in C(K;\R)$ 和 $\varepsilon>0$。对任意 $x\ne y$，由分离点性
取 $h\in A$ 使 $h(x)\ne h(y)$。对 $h$ 作仿射变换，可得
$g_{x,y}\in A$ 满足
$g_{x,y}(x)=f(x)$、$g_{x,y}(y)=f(y)$；$x=y$ 时取常函数 $f(x)$。

先固定 $x$。对每个 $y$，连续性给出 $y$ 的邻域 $V_y$，其上
$g_{x,y}>f-\varepsilon$。有限个 $V_{y_1},\ldots,V_{y_r}$ 覆盖 $K$。由
引理 \ref{lem:stone-weierstrass-lattice-closure}，
\[
 g_x=g_{x,y_1}\vee\cdots\vee g_{x,y_r}\in\overline A.
\]
于是 $g_x>f-\varepsilon$ 于 $K$，且 $g_x(x)=f(x)$。再由连续性，存在
$x$ 的邻域 $U_x$，其上 $g_x<f+\varepsilon$。

有限个 $U_{x_1},\ldots,U_{x_s}$ 覆盖 $K$。令
$g=g_{x_1}\wedge\cdots\wedge g_{x_s}\in\overline A$。每个 $g_{x_i}$ 都
大于 $f-\varepsilon$，而每一点落在某个 $U_{x_i}$ 中，故
\[
 f-\varepsilon<g<f+\varepsilon\quad\text{于 }K.
\]
于是 $\norm{g-f}_\infty\le\varepsilon$。因为 $g\in\overline A$，还可从
$A$ 中选函数任意逼近 $g$；合并误差即知 $f\in\overline A$。
\end{proof}

对复函数，分离点和含常数还不够。需要 $A$ 对复共轭封闭，即
$f\in A\Rightarrow\bar f\in A$。令 $A_{\R}$ 为 $A$ 中的实值函数。共轭
封闭保证 $\operatorname{Re}f,\operatorname{Im}f\in A_{\R}$；若
$f(x)\ne f(y)$，二者中至少一个分离 $x,y$，所以 $A_{\R}$ 仍分离点。
对实部和虚部分别应用实定理便得到：

\begin{corollary}[复 Stone--Weierstrass 定理]
\label{cor:complex-stone-weierstrass}
设 $A\subset C(K;\C)$ 是含常数、分离点且对共轭封闭的复子代数，则
$\overline A=C(K;\C)$。
\end{corollary}

几个稠密性结论现在可由同一判据得到：
\begin{enumerate}
 \item $K\subset\R^n$ 紧时，坐标多项式限制在 $K$ 上含常数并分离点，故
 多元多项式在 $C(K;\R)$ 中稠密；
 \item 圆周上的三角多项式由 $e^{int}$（$n\in\Z$）生成，含共轭并分离点，
 故在连续周期函数中一致稠密；
 \item 若复代数不含共轭，结论可能失败。在单位圆
 $\mathbb T=\{z:|z|=1\}$ 上，只含非负幂的多项式构成含常数并分离点的
 复代数，却不能一致逼近 $\bar z$。事实上，连续线性泛函
 \[
  L(f)=\int_0^{2\pi}f(e^{it})e^{it}\,dt
 \]
 在每个多项式 $\sum_{k=0}^Na_kz^k$ 上取零，而 $L(\bar z)=2\pi$；故
 $\bar z$ 不属于该代数的一致闭包。
\end{enumerate}

\section*{习题}
\addcontentsline{toc}{section}{习题}

\subsection*{A组　Bernstein 构造}
\begin{enumerate}[label=\textbf{\thechapter.\arabic*.}]
 \exerciseitem{bernstein-moment-identities}
 从二项式定理完整推导三个 Bernstein 矩恒等式。
 \exerciseitem{bernstein-positive-contraction}
 证明 $B_n$ 是正线性算子、保持常数，并满足
 $\norm{B_nf}\le\norm f$。
 \exerciseitem{bernstein-modulus-estimate}
 补全定理 \ref{thm:bernstein-uniform-approximation} 的每一步估计。
 \exerciseitem{bernstein-lipschitz-rate}
 若 $f$ 的 Lipschitz 常数为 $L$，证明
 $\norm{B_nf-f}_\infty\le L/(2\sqrt n)$。
 \exerciseitem{bernstein-linear-quadratic}
 计算 $B_n1$、$B_nt$、$B_nt^2$，并判断哪些多项式被精确保持。
 \exerciseitem{bernstein-monotonicity}
 用相邻差分证明 $f$ 单调时 $B_nf$ 单调。
\end{enumerate}

\subsection*{B组　稠密性与格方法}
\begin{enumerate}[label=\textbf{\thechapter.\arabic*.},resume]
 \exerciseitem{fixed-degree-polynomials-closed}
 证明次数不超过 $N$ 的多项式在 $C([a,b])$ 中构成闭子空间。
 \exerciseitem{dense-subspace-propagates-linear-identity}
 证明连续线性泛函若在稠密子集上为零，则在全空间为零。
 \exerciseitem{absolute-value-polynomial-approximation}
 说明引理 \ref{lem:stone-weierstrass-lattice-closure} 中可要求逼近多项式
 满足 $p_n(0)=0$，并完成复合函数的一致估计。
 \exerciseitem{two-point-interpolation-from-separation}
 证明分离点的含常数代数能为任意 $x\ne y$ 和任意实数 $a,b$ 构造
 $g(x)=a,g(y)=b$ 的函数 $g$。
 \exerciseitem{stone-weierstrass-finite-cover-proof}
 重写 Stone--Weierstrass 证明，逐一说明两次有限子覆盖和最大、最小运算
 的作用。
\end{enumerate}

\subsection*{C组　典型稠密性结论}
\begin{enumerate}[label=\textbf{\thechapter.\arabic*.},resume]
 \exerciseitem{multivariate-polynomial-density}
 对任意紧集 $K\subset\R^n$，证明多元实多项式在 $C(K;\R)$ 中稠密。
 \exerciseitem{trigonometric-polynomial-density}
 用复 Stone--Weierstrass 定理证明三角多项式在连续 $2\pi$-周期函数中一致
 稠密。
 \exerciseitem{even-polynomial-density}
 判断偶多项式在 $C([-1,1])$ 中的闭包，并证明你的结论。
 \exerciseitem{polynomials-with-zero-at-origin}
 求满足 $p(0)=0$ 的多项式在 $C([0,1])$ 中的闭包。
 \exerciseitem{complex-conjugation-necessity}
 设 $\mathbb T=\{z\in\C:|z|=1\}$。证明只含非负频率的多项式
 $\sum_{k=0}^Na_kz^k$ 不能在 $\mathbb T$ 上一致逼近 $\bar z$。
 \optionalexerciseitem{bernstein-convexity-preservation}
 用二阶差分证明 Bernstein 算子保持凸性，并由此说明它对凸函数给出何种
 单侧估计。
\end{enumerate}
